\usection{Técnicas e Instrumentos de Recolección de Datos}

\citeA{arias2012} define las técnicas de recolección de datos como “el procedimiento o forma particular de obtener los datos o información” (p.~67), mientras que los instrumentos de recolección de datos son “cualquier recurso, dispositivo o formato (en papel o digital), que se pueda utilizar para obtener, registrar o almacenar información” (p.~68).

En este sentido, la recolección de datos en la fase de análisis del proyecto consistió en el uso de la revisión documental, que implicó la búsqueda, lectura y análisis de fuentes bibliográficas relacionadas con las distintas dimensiones del estudio, tales como el pensamiento computacional, el currículo educativo en Venezuela, la concurrencia y el desarrollo de lenguajes de programación; incluyendo entre los tipos de fuentes libros, artículos académicos y sitios web reconocidos. Esta técnica permitió obtener información relevante para fundamentar teóricamente el proyecto.
